\chapter[Resultados Experimentais Completos]{Resultados Experimentais Completos}
\label{ape:resultados}

Este apêndice apresenta as tabelas completas de resultados experimentais para todas as 30 instâncias e métodos avaliados. As tabelas do corpo principal (Capítulo~\ref{experimentos}) apresentam versões resumidas ou agregadas; as tabelas abaixo incluem os valores por instância.

\section{Cobertura Experimental}

A \autoref{tab:ape-cobertura} resume a cobertura experimental completa.

\begin{table}[htbp]
\centering
\caption{Cobertura experimental completa.}
\label{tab:ape-cobertura}
% Generated by scripts/gerar-metricas-monografia.py; do not edit manually.
% Source: src/data/results via scripts/gerar-metricas-monografia.py.
\begin{tabular}{lrrr}
\toprule
Metodo & Runs & Instancias & Sementes \\
\midrule
ACO & 1530 & 30 & 51 \\
GA & 1530 & 30 & 51 \\
PSO & 1530 & 30 & 51 \\
Lower bound & 30 & 30 & 1 \\
Busca exaustiva & 18 & 18 & 1 \\
\bottomrule
\end{tabular}

\legend{Fonte: dados da pesquisa.}
\end{table}

\section{\'Otimes da Busca Exaustiva}

A \autoref{tab:ape-otimos-bf} apresenta os \'otimos encontrados pela busca exaustiva para as instâncias de 10 a 15 pontos.

\begin{table}[htbp]
\centering
\caption{\'Otimes da busca exaustiva.}
\label{tab:ape-otimos-bf}
% Generated by scripts/gerar-metricas-monografia.py; do not edit manually.
% Source: src/data/results via scripts/gerar-metricas-monografia.py.
\begin{tabular}{rrrrrr}
\toprule
Instancia & Otimo & Instancia & Otimo & Instancia & Otimo \\
\midrule
10a & 8{,}251 & 10b & 7{,}906 & 10c & 9{,}545 \\
11a & 9{,}493 & 11b & 8{,}823 & 11c & 10{,}260 \\
12a & 9{,}501 & 12b & 9{,}769 & 12c & 10{,}480 \\
13a & 10{,}950 & 13b & 11{,}474 & 13c & 10{,}042 \\
14a & 11{,}557 & 14b & 11{,}467 & 14c & 11{,}486 \\
15a & 9{,}129 & 15b & 12{,}185 & 15c & 12{,}160 \\
\bottomrule
\end{tabular}

\legend{Fonte: dados da pesquisa.}
\end{table}

\section{Gap vs.\ Busca Exaustiva}

A \autoref{tab:ape-gap-bf} apresenta o gap percentual de cada m\'etodo em rela\c{c}\~ao ao \'otimo nas 18 instâncias com \'otimo conhecido.

\begin{table}[htbp]
\centering
\caption{Gap percentual em rela\c{c}\~ao ao \'otimo.}
\label{tab:ape-gap-bf}
% Generated by scripts/gerar-metricas-monografia.py; do not edit manually.
% Source: src/data/results via scripts/gerar-metricas-monografia.py.
\begin{tabular}{lrrrrr}
\toprule
Metodo & Runs & Gap medio & Gap min. & Gap max. & Taxa de acerto \\
\midrule
ACO & 918 & 0{,}4296\% & 0{,}0000\% & 9{,}4445\% & 71{,}79\% \\
GA & 918 & 5{,}2814\% & 0{,}0000\% & 41{,}4267\% & 26{,}47\% \\
PSO & 918 & 22{,}2915\% & 0{,}0000\% & 76{,}5549\% & 4{,}25\% \\
\bottomrule
\end{tabular}

\legend{Fonte: dados da pesquisa.}
\end{table}

\section{Gap vs.\ Lower Bound AP}

A \autoref{tab:ape-lb-gap} apresenta o gap percentual de cada m\'etodo em rela\c{c}\~ao ao lower bound AP para todas as 30 instâncias.

\begin{table}[htbp]
\centering
\caption{Gap percentual em rela\c{c}\~ao ao lower bound AP.}
\label{tab:ape-lb-gap}
% Generated by scripts/gerar-metricas-monografia.py; do not edit manually.
% Source: src/data/results via scripts/gerar-metricas-monografia.py.
\begin{tabular}{lr}
\toprule
Metrica & Valor \\
\midrule
Gap medio LB$\to$BF & 51{,}36\% \\
Gap minimo & 40{,}79\% \\
Gap maximo & 65{,}35\% \\
Instancias avaliadas & 18 \\
Violacoes LB $\leq$ BF & 0 \\
\bottomrule
\end{tabular}

\legend{Fonte: dados da pesquisa.}
\end{table}

\section{Resultados em Instâncias Grandes}

A \autoref{tab:ape-grandes} apresenta os resultados agregados para as instâncias de 50 e 100 pontos.

\begin{table}[htbp]
\centering
\caption{Makespan nas instâncias grandes.}
\label{tab:ape-grandes}
% Generated by scripts/gerar-metricas-monografia.py; do not edit manually.
% Source: src/data/results via scripts/gerar-metricas-monografia.py.
\begin{tabular}{lrrrrrr}
\toprule
Instancia & ACO best & ACO media & GA best & GA media & PSO best & PSO media \\
\midrule
50a & 24{,}6688 & 25{,}6475 & 41{,}0120 & 45{,}3925 & 53{,}7651 & 59{,}7284 \\
50b & 26{,}3528 & 27{,}4843 & 42{,}6487 & 48{,}6345 & 57{,}0330 & 64{,}3194 \\
50c & 26{,}2372 & 27{,}0062 & 42{,}0122 & 46{,}9438 & 54{,}4114 & 62{,}3455 \\
100a & 42{,}6718 & 45{,}3528 & 97{,}5921 & 110{,}6608 & 127{,}4034 & 136{,}5131 \\
100b & 45{,}4689 & 48{,}0596 & 97{,}5343 & 110{,}4959 & 126{,}3967 & 135{,}9987 \\
100c & 45{,}4858 & 46{,}9487 & 95{,}1632 & 111{,}3243 & 126{,}1838 & 138{,}1908 \\
\bottomrule
\end{tabular}

\legend{Fonte: dados da pesquisa.}
\end{table}

\section{Tempo Computacional}

A \autoref{tab:ape-tempo-100} apresenta o tempo de otimiza\c{c}\~ao para as instâncias de 100 pontos.

\begin{table}[htbp]
\centering
\caption{Tempo de otimiza\c{c}\~ao nas instâncias de 100 pontos.}
\label{tab:ape-tempo-100}
% Generated by scripts/gerar-metricas-monografia.py; do not edit manually.
% Source: src/data/results via scripts/gerar-metricas-monografia.py.
\begin{tabular}{lrrrrr}
\toprule
Instancia & ACO ms & GA ms & PSO ms & ACO/GA & ACO/PSO \\
\midrule
100a & 4262{,}06 & 38{,}82 & 76{,}24 & 109{,}78x & 55{,}91x \\
100b & 4276{,}55 & 36{,}27 & 75{,}25 & 117{,}89x & 56{,}83x \\
100c & 4266{,}55 & 39{,}37 & 73{,}84 & 108{,}36x & 57{,}78x \\
\bottomrule
\end{tabular}

\legend{Fonte: dados da pesquisa.}
\end{table}
